\section{Существующие подходы к анализу связи между нейронной активностью и поведением}
\label{sec:behavior_encoding}

Выявление связей между активностью отдельных нейронов и поведением животного является одной из центральных задач нейронауки.
На популяционном уровне нейронные сигналы формируют высокоразмерные представления, направляющие поведение и когнитивные процессы~\cite{Averbeck2006}.
Однако при регистрации активности у свободно поведающих животных основную трудность представляет не само обнаружение связи между нейронной активностью и поведением, а корректная атрибуция: какая именно поведенческая переменная объясняет активность данного нейрона, если поведение многомерно, его компоненты коррелированы между собой и имеют сложную временн\'{у}ю структуру~\cite{Kriegeskorte2019}.

Проблема атрибуции особенно обостряется в свободном поведении, поскольку поведенческие переменные связаны друг с другом.
Современные средства автоматического анализа поведения --- как неконтролируемые (MoSeq)~\cite{Weinreb2024}, так и контролируемые (DeepAction~\cite{Harris2023deepaction}, MARS~\cite{Segalin2021}) и основанные на правилах (BehaviorDEPOT)~\cite{Gabriel2022} --- позволяют извлекать из видеозаписей десятки и сотни поведенческих признаков. Результатом становятся насыщенные, но высоко коррелированные наборы данных, порождающие проблему идентифицируемости: видимая селективность нейрона к одной переменной может быть следствием его настройки на коррелированный ковариат.
Например, нейрон, на первый взгляд кодирующий локомоцию, на самом деле может быть настроен на скорость, а клетка места --- ошибочно интерпретирована как реагирующая на социальное взаимодействие лишь потому, что контакты приурочены к определённым местам. Эта проблема широко распространена в нейронауке: непроинструктированные движения объясняют значительную долю активности коры~\cite{Musall2019}, а тета-модуляция затрудняет разделение кодирования скорости и фазы в гиппокампе~\cite{Vanderwolf1969, McNaughton1983}.

Существующие подходы к анализу связи между нейронной активностью и поведением можно разделить на три группы в зависимости от уровня анализа и типа получаемой информации.


\subsection{Подходы на основе качества кодирующих моделей}
\label{sec:decoding_approaches}

Первая группа подходов оценивает, насколько хорошо поведение предсказывается по нейронной активности с помощью моделей машинного обучения --- линейных и нелинейных декодеров, вероятностных моделей и глубоких нейронных сетей. Этот подход позволяет количественно оценить информационную ёмкость нейронного представления: если декодер точно воспроизводит траекторию движения, положение головы или дискретные поведенческие состояния, значит популяция нейронов содержит соответствующую информацию~\cite{DayanAbbott2001}.

Однако высокое качество декодирования не обязательно означает понимание кодирующего механизма.
Вероятностные декодеры способны предсказывать сложное поведение, но при этом не дают ответа на вопрос, какие именно поведенческие переменные объясняют активность конкретного нейрона --- особенно когда ковариаты коррелированы~\cite{Kriegeskorte2019}.
Линейные меры, в свою очередь, не способны обнаруживать нелинейные зависимости, что ограничивает полноту анализа. В результате декодирующие модели характеризуют популяционное представление в целом, но не позволяют определить вклад отдельных нейронов и тип кодируемой ими информации.


\subsection{Подходы на основе анализа популяционного кода}
\label{sec:population_approaches}

Второй класс методов направлен на выявление скрытой структуры нейронной активности на уровне популяции. Методы снижения размерности --- анализ главных компонент (PCA), t-SNE, UMAP и другие --- проецируют многомерную нейронную активность в пространство малой размерности, позволяя визуализировать коллективную динамику и выявлять групповые паттерны~\cite{Cunningham2014}. Такой подход оказался продуктивным для обнаружения нейронных многообразий и описания траекторий популяционной активности, однако при редукции размерности вклад индивидуальных нейронов в наблюдаемые паттерны оказывается скрыт, что не позволяет определить, какие нейроны и к каким переменным селективны.

Информационно-теоретические подходы на уровне популяции количественно оценивают информацию, передаваемую совокупностью нейронов. В частности, недавно разработанный инструмент MINT~\cite{Lorenz2025} специально предназначен для оценки популяционной передачи информации. Однако и эти методы ориентированы на ансамблевые характеристики и не позволяют установить, какой конкретный нейрон связан с какой конкретной переменной.

Таким образом, популяционные подходы дают общую картину информационной структуры нейронной активности, но не решают задачу атрибуции на уровне отдельных клеток.


\subsection{Подходы на основе индивидуальной селективности}
\label{sec:selectivity_approaches}

Третий подход направлен на определение селективности каждого нейрона к конкретным внешним переменным.
Именно этот уровень анализа позволяет выявлять функциональные типы клеток --- клетки места, клетки направления головы, нейроны скорости и т.\,п. --- и устанавливать, как кодирующие свойства отдельных нейронов формируют популяционное представление в целом.

Информационно-теоретические методы, основанные на взаимной информации (MI), предоставляют для этого принципиальную основу, поскольку MI улавливает произвольные статистические зависимости --- как линейные, так и нелинейные --- и равна нулю тогда и только тогда, когда переменные независимы.

Тем не менее существующие инструменты информационно-теоретического анализа ориентированы на различные экспериментальные парадигмы и уровни анализа и в отдельности не решают задачу индивидуальной селективности в контексте непрерывного кальциевого имиджинга:

\begin{itemize}
    \item \textbf{NIT (Neuroscience Information Toolbox)}~\cite{Maffulli2022} --- наиболее полный инструмент для оценки MI с коррекцией смещения и подробными рекомендациями по работе с кальциевыми данными. Однако NIT разработан для парадигмы повторных проб (trial-based), предполагающей многократное предъявление стимулов в контролируемых условиях, и не приспособлен к непрерывным записям свободного поведения, где необходимо учитывать временн\'{у}ю автокорреляцию сигналов, оптимизировать задержки между нейронной активностью и поведением и контролировать ковариацию поведенческих переменных.

    \item \textbf{MINT}~\cite{Lorenz2025} --- инструмент для оценки передачи информации на популяционном уровне, не предназначенный для индивидуального анализа отдельных нейронов.

    \item \textbf{FRITES}~\cite{Combrisson2022} --- фреймворк группового статистического вывода для электрофизиологических данных, ориентированный на сопоставление групп испытуемых, а не на характеризацию селективности отдельных клеток в кальциевых записях.

    \item \textbf{IDTxl}~\cite{IDTxl2019} --- инструмент для анализа сетевой связности (transfer entropy, причинность по Грейнджеру), направленный на выявление информационных потоков между узлами сети, а не на оценку селективности отдельных нейронов к внешним переменным.

    \item \textbf{HOI}~\cite{Neri2024} --- инструмент для анализа взаимодействий высокого порядка между переменными, дополняющий попарные методы, но не ориентированный на задачу индивидуальной селективности.

    \item \textbf{GCMI}~\cite{Ince2017} --- эффективный метод оценки MI через гауссову копулу, работающий с непрерывными данными. Однако оценка MI сама по себе не решает ключевых задач анализа свободного поведения: выбора оптимальных временных задержек при наличии автокорреляции и разделения истинного кодирования от артефактов, порождённых ковариацией поведенческих переменных.
\end{itemize}

Помимо информационно-теоретических инструментов, кальциевый имиджинг предъявляет к методам анализа дополнительные требования.
Кальциевые индикаторы, включая наиболее быстрые (например, jGCaMP8), характеризуются временами затухания от десятков до сотен миллисекунд~\cite{Zhang2023}, что приводит к временн\'{о}й свёртке сигнала: флуоресценция от предшествующих событий ложно коррелирует с более поздними поведенческими эпизодами~\cite{Wei2020}.
В отличие от спайковых данных, кальциевый сигнал непрерывен и не может быть описан точечным процессом, а деконволюция вносит модельно-зависимые искажения и может снижать качество последующего декодирования или атрибуции~\cite{Rupprecht2021}.
При этом регистрация охватывает от сотен до тысяч нейронов одновременно~\cite{Chen2013, Zong2022}, а поведение описывается переменными принципиально разных типов --- от дискретных категориальных состояний (груминг, стойка на задних лапах) до непрерывных величин (скорость, направление головы), что требует аналитических подходов, способных работать с гетерогенными типами данных в рамках единого фреймворка.

Таким образом, в существующей литературе отсутствует инструмент, который одновременно решал бы следующие задачи: работа непосредственно с непрерывным кальциевым сигналом без предварительной деконволюции; обработка смешанных типов поведенческих переменных; оптимизация временных задержек с учётом кинетики индикатора и проспективного/ретроспективного кодирования; разделение истинной селективности от артефактов, порождённых ковариацией поведенческих признаков, --- и всё это в масштабе тысяч нейронов и десятков поведенческих переменных. Заполнению этого пробела посвящена Глава~\ref{sec:individual}, где описан фреймворк INTENSE.
